\chapter{Catalog of Defects and Dead Code in the Reference Implementation}
\label{ch:defects}

Every entry below was found during an exhaustive, function-by-function cross-reference of the
reference C++ source against an independent re-implementation, and confirmed either by
constructing a triggering test case or by an exhaustive grep-based call-site audit (for
dead-code claims: a claim of ``zero call sites'' means every reference to the symbol in the
entire source tree was individually inspected). None of these were assumed from documentation or
comments alone.

\section{Summary table}

\begin{center}
\begin{tabular}{@{}p{0.6cm}p{3.6cm}p{2.2cm}p{2.1cm}p{5.2cm}@{}}
\toprule
\# & \textbf{Defect} & \textbf{Affects live output?} & \textbf{Severity} & \textbf{Recommended
  handling} \\
\midrule
D1 & Uninitialized LUT arrays & Only non-ACGT input & Latent UB & Zero-initialize (\S
  \ref{sec:d1}) \\
D2 & \code{kmers\_notmatched} miscount & Diagnostic only & Cosmetic & Reorder the increment (\S
  \ref{sec:d2}) \\
D3 & \code{more\_seeds\_if\_cheap} unreachable & No (dead) & None & Optional to wire in (\S
  \ref{sec:d3}) \\
D4 & Hardcoded 100-segment cap & Crashes on $>100$ contigs & Usability & Size dynamically (\S
  \ref{sec:d4}) \\
D5 & \code{Buckets::size()} dead stub & No (dead) & None & Drop (\S\ref{sec:d5}) \\
D6 & \code{Mapping::do\_overlap} dead & No (dead) & None & Drop (\S\ref{sec:d6}) \\
D7 & \code{Matcher::do\_overlap}/\newline\code{lost\_correct\_mapping} dead & No (dead) & None &
  Drop (\S\ref{sec:d7}) \\
D8 & Dead \code{lmax} parameter in ground-truth \code{bestFixedLength} & Diagnostic only &
  Cosmetic/confusing & Drop the parameter (\S\ref{sec:d8}) \\
D9 & Off-by-one in \code{bestIncludedJaccard} & Diagnostic only (\code{-v 2}) & \textbf{Real
  correctness bug} & Fix indexing (\S\ref{sec:d9}) \\
D10 & Missing bounds clamp in ground-truth \code{collect\_matches} & Diagnostic only, potential
  crash/OOB & Latent bug & Clamp to segment length (\S\ref{sec:d10}) \\
D11 & \code{RefSegment::seq} unused waste & No (memory only) & Efficiency & Omit field (\S
  \ref{sec:d11}) \\
D12 & Dead \code{-a}/SAM output path & No (dead) & None & Omit entirely (\S\ref{sec:d12}) \\
D13 & Dead \code{-S}/\code{-n} CLI flags & No (dead) & None & Keep parsed-but-inert for
  compatibility (\S\ref{sec:d13}) \\
D14 & \code{one\_sweep}/\code{-B} never branched on & No (dead) & Usability/API honesty &
  Implement a concrete meaning (\S\ref{sec:d14}) \\
D15 & \textbf{Per-read counters never reset} & \textbf{Yes --- live PAF tags} & \textbf{Real
  correctness bug} & Fully reset every read (\S\ref{sec:d15}) \\
D16 & \textbf{Unconditional dereference of empty optional} & \textbf{Yes --- crashes/UB on every
  unmapped read} & \textbf{Real correctness bug (UB)} & Emit a defined minimal record (\S
  \ref{sec:d16}) \\
D17 & \code{gt\_C\_r\_lmax} never assigned & Diagnostic only (\code{-v 2}) & Cosmetic & Leave
  unassigned or compute it properly (\S\ref{sec:d17}) \\
D18 & PAF columns 10/11 always repeat query length & \textbf{Yes, if spec-conformance matters} &
  Cosmetic/compat & Populate real values if spec conformance is required (\S\ref{sec:d18}) \\
D19 & \code{Timers::secs}'s try/catch never catches anything & No (silent no-op) & Latent & Handle
  missing keys directly (\S\ref{sec:d19}) \\
D20 & \code{ParsedQueryId::parse} catches then rethrows & No & Style only & Simplify (\S
  \ref{sec:d20}) \\
\bottomrule
\end{tabular}
\end{center}

\section{D1: Uninitialized hash lookup tables}
\label{sec:d1}
\textbf{Location:} \code{FracMinHash} construction, \texttt{src/sketch.h} (\S\ref{sec:sketch-lut}).
\textbf{Description:} \code{LUT\_fw}/\code{LUT\_rc} are fixed 256-entry arrays; only the 8 entries
for \texttt{ACGTacgt} are ever written. Any other input byte (most commonly \texttt{N}) reads
uninitialized memory when used as a hash contribution --- undefined behavior in C++, and not even
a fixed garbage value (it depends on prior stack contents at the time \code{sketch()} first runs).
\textbf{Trigger:} any reference or read sequence containing a non-ACGT byte.
\textbf{Fix:} zero-initialize both tables before filling in the 8 real entries. This is a strict
safety improvement with zero behavioral change on ACGT-only input (the overwhelming majority of
real data), and was adopted without further discussion in the companion re-implementation.

\section{D2: \code{kmers\_notmatched} always reports the full sketch size}
\label{sec:d2}
\textbf{Location:} \code{unique\_elements\_with\_info}, \texttt{src/shmap.h} (\S
\ref{sec:seeding-grouping}). \textbf{Description:} the increment
\code{if (hits\_in\_t > 0) nonzero += strike;} is placed textually after \code{strike} has
already been reset to 0 for the next group, so it always adds zero; \code{kmers\_notmatched} (a
diagnostic counter only) therefore always reports the entire sketch size regardless of actual
match rate. \textbf{Impact:} cosmetic only --- feeds one end-of-run percentage in the stderr
summary (\S\ref{sec:runtime-warnings}), nothing else. \textbf{Fix:} move the increment before the
reset (or reorder the two statements).

\section{D3: \code{more\_seeds\_if\_cheap} is fully implemented but unreachable}
\label{sec:d3}
\textbf{Location:} \texttt{src/shmap.h} (\S\ref{sec:seed-budget}). \textbf{Description:} a
complete, presumably-tested function to opportunistically extend the seed budget when the next
seeds are (almost) unique in the reference; its only call site in \code{map\_read} is commented
out. \textbf{Impact:} none on shipped behavior (never runs). \textbf{Recommendation:} a
from-scratch implementation may wire it in (a plausible, low-risk recall improvement) or omit it
to match shipped behavior exactly; either is defensible.

\section{D4: Hardcoded 100-segment cap}
\label{sec:d4}
\textbf{Location:} \code{Buckets} construction, \texttt{src/buckets.h} (\S\ref{sec:buckets-struct}).
\textbf{Description:} \code{MAX\_SEGMENTS = 100}; construction throws for references with more
segments. \textbf{Trigger:} draft/scaffolded assemblies with hundreds or thousands of small
contigs --- not a rare input in practice. \textbf{Fix:} size the per-segment bucket storage to the
reference's actual segment count.

\section{D5--D7: Dead stubs and unreachable methods}
\label{sec:d5}
\label{sec:d6}
\label{sec:d7}
Three separate pieces of code are fully defined but have \textbf{zero} call sites anywhere in the
source tree (verified exhaustively, not just ``no obvious caller''):
\begin{itemize}
  \item \code{Buckets::size()} (\texttt{src/buckets.h}) --- literally \code{return -1;} with a
    \code{// TODO: implement} comment.
  \item \code{Mapping::do\_overlap} (\texttt{src/sketch.h}) --- a boolean overlap check,
    superseded everywhere by the (used) \code{Mapping::overlap} that returns a fraction
    (\S\ref{sec:overlap-exclusion}).
  \item \code{Matcher::do\_overlap} and \code{Matcher::lost\_correct\_mapping}
    (\texttt{src/refine.h}) --- ground-truth-overlap-checking helpers, superseded by the
    \code{-v 2} diagnostic path (\S\ref{sec:ground-truth}) actually used.
\end{itemize}
\textbf{Recommendation:} omit all three from a from-scratch implementation; they contribute no
behavior.

\section{D8: Dead \code{lmax} parameter in the ground-truth \code{bestFixedLength}}
\label{sec:d8}
\textbf{Location:} \texttt{src/refine.h}, the \code{Matcher}-scoped \code{bestFixedLength}
(\S\ref{sec:matcher-best-fixed-length}). \textbf{Description:} the function signature accepts an
\code{lmax} parameter, and every call site dutifully computes and passes a value for it
(\code{2*bucket\_l} or \code{bucket\_l}), but the function body's actual window-width check uses
\code{P\_sz} instead --- a commented-out \code{lmax}-based condition sits directly above the live
one, evidence the author switched approaches and left the now-ornamental parameter (and its
call-site computation) in place. \textbf{Consequence:} the two ground-truth tags
\code{gt\_C\_l\_lmax} and \code{gt\_C\_l} (\S\ref{sec:ground-truth}), whose call sites differ only
in what they pass for this dead parameter, are numerically \emph{identical} in practice.
\textbf{Fix:} drop the parameter; it does nothing.

\section{D9: Off-by-one in \code{bestIncludedJaccard} (real bug)}
\label{sec:d9}
\textbf{Location:} \texttt{src/refine.h}, \code{Matcher::bestIncludedJaccard}
(\S\ref{sec:matcher-best-fixed-length}). Full technical description already given in
\S\ref{sec:defect-prevr}; repeated here for completeness of this catalog. \textbf{Description:}
the function's window-growing loop scores mid-iteration using the position \emph{one before} the
just-processed element (\code{prev(r)}) for the reported span, while \code{intersection} at that
point already reflects the just-processed element --- an internal inconsistency that can and does
produce \code{intersection > span}, violating the metric's own $[0,1]$-boundedness invariant. The
original source carries an explicit \code{// TODO: should it be prev(r) instead?} comment on this
exact line. \textbf{Trigger:} confirmed via a small constructed example (a $\sim$50bp read against
a $\sim$200bp synthetic reference, exercising the \code{-v 2} ground-truth ``included Jaccard''
computation). \textbf{Impact:} affects only the \code{-v 2} diagnostic tags/TSV, never the primary
PAF output or any mapping decision (this sweep is not used by the main scoring path,
\S\ref{sec:best-fixed-length}, which has the analogous computation \emph{and gets it right} by
scoring only after its inner loop fully exits). \textbf{Fix:} score using the \emph{current},
just-processed right-boundary element consistently (matching what \code{intersection} already
counts), not the one before it.

\section{D10: Missing bounds clamp in ground-truth \code{collect\_matches}}
\label{sec:d10}
\textbf{Location:} \texttt{src/refine.h}, \code{Matcher::collect\_matches}. \textbf{Description:}
this function computes a bucket's \code{[start, end)} range into a segment's sketch array but does
not clamp \code{end} to the segment's actual length --- for the last bucket of a segment, a
bucket's nominal end can run past the sketch array's bounds, risking an out-of-bounds read. The
same file's own (separately dead, D7) \code{do\_overlap} function \emph{does} apply exactly this
clamp for the same quantity, suggesting the omission here is an oversight rather than intentional.
\textbf{Fix:} clamp \code{end} to \code{min(end, len(segment.kmers))} before iterating, mirroring
the clamp already present elsewhere in the same file.

\section{D11: \code{RefSegment::seq} is stored but never read (in the shipped feature set)}
\label{sec:d11}
\textbf{Location:} \texttt{src/sketch.h} (\S\ref{sec:type-refsegment}). \textbf{Description:} the
full raw nucleotide sequence of every reference segment is retained in memory, but the only code
that ever reads it (a SAM-output/edit-distance feature) is entirely commented out and
unreachable. For large references this is a substantial, unnecessary memory cost (roughly doubling
index memory beyond the sketch itself). \textbf{Recommendation:} omit the field entirely in a
from-scratch implementation targeting the mapping functionality this document describes; only
retain it if also implementing base-level alignment output.

\section{D12: Dead \code{-a}/SAM output path}
\label{sec:d12}
\textbf{Location:} CLI parsing (\texttt{src/io.h}) and \texttt{src/sketch.h}.
\textbf{Description:} the \code{-a} flag (\code{sam}) is parsed and stored but never read by any
live code; the SAM-format output function it would have selected is fully commented out (and
itself depends on an edit-distance library, \code{edlib}, that is otherwise entirely unused by the
shipped binary). \textbf{Recommendation:} omit this flag from a from-scratch implementation rather
than reproduce a no-op switch.

\section{D13: Dead \code{-S}/\code{-n} CLI flags}
\label{sec:d13}
\textbf{Location:} CLI parsing (\texttt{src/io.h}). \textbf{Description:} \code{max\_seeds}
(\code{-S}) and \code{normalize} (\code{-n}) are parsed, range-validated, and included in the
parameter dump, but neither is ever read by any algorithmic code. \textbf{Distinction from D12:}
unlike \code{-a}, these are not accompanied by a larger commented-out feature --- they appear to be
either forward-looking placeholders or remnants of a removed feature. \textbf{Recommendation:}
unlike the confirmed-and-explicitly-flagged \code{-a} case, a from-scratch implementation should
by default \emph{keep} these as accepted-but-inert CLI parameters (for command-line compatibility
with existing scripts/pipelines that may already pass them) rather than unilaterally removing
CLI-visible surface area that wasn't specifically identified for removal.

\section{D14: \code{one\_sweep}/\code{-B} has no implemented behavior}
\label{sec:d14}
\textbf{Location:} template parameter and CLI flag threaded throughout \texttt{src/shmap.h},
\texttt{src/mapper.cpp}. \textbf{Description:} unlike \code{no\_bucket\_pruning}/\code{-P} (which
does gate real behavior, \S\ref{sec:seed-heuristic-pass}) or \code{abs\_pos}/\code{-F} (which
gates coordinate-space arithmetic throughout Chapters~\ref{ch:bucketing}--\ref{ch:scoring}),
\code{one\_sweep} is threaded through every relevant function signature and even given a
descriptive CLI help string (``disregards seed heuristic and runs one sweep on all matches''), but
is never once tested inside any conditional in the mapper's actual logic --- confirmed via
exhaustive grep. \textbf{Recommendation:} since this document's companion implementation was
tasked with making every configuration combination genuinely meaningful (rather than silently
inert), it defines a concrete best-effort semantics matching the flag's own help text: consume
\emph{every} seed (\code{S := len(seeds)}) rather than the theta-derived budget
(\S\ref{sec:seed-budget}), leaving the rest of the pipeline (bucket pruning, refinement)
unchanged. This is an explicit design choice filling a gap the reference source left open, not a
recovered original behavior --- a from-scratch implementation is free to define \code{one\_sweep}
differently (e.g.\ literally skip bucket-based clustering and score one global sweep over every
raw match, which is closer to a literal reading of the help text but a substantially larger
undertaking) as long as it documents which interpretation it chose.

\section{D15: Per-read counters are never actually reset (real bug, affects live output)}
\label{sec:d15}
\textbf{Location:} \code{map\_read}, \texttt{src/shmap.h} (\S\ref{sec:map-read-counters}).
\textbf{Description:} the call that would fully reset the per-read \code{Counters} object
(\code{C.clear()}) is commented out; the subsequent re-initialization only names a subset of the
counter keys the function goes on to increment. The following counters are left unreset across
every read a single mapper instance processes: \code{read\_len}, \code{total\_matches},
\code{possible\_matches}, \code{kmers\_sketched}, \code{kmers\_unique}, \code{kmers\_seeds},
\code{seeded\_buckets}. Because the (also unreset) running counters object is merged into the
run-wide totals \emph{after every single read}, these specific counters compound
super-linearly across a run. \textbf{Concretely affected live output:} \code{total\_matches:i:}
and its derived \code{match\_inefficiency:f:} PAF tag (\S\ref{ch:output}) are read \emph{directly}
from this same, never-reset counters object \emph{before} the end-of-read merge into run-wide
totals --- meaning every mapped read's own reported \code{total\_matches:i:} value already
reflects leftover accumulation from every previous read the same mapper instance processed, not
just the current read. This is the most consequential defect in this catalog: it corrupts an
observable, per-mapping output field, not merely an end-of-run summary statistic.
\textbf{Confirmed} by observing, in a controlled two-read test, \code{total\_matches:i:} values
that were consistent per-read only after applying the fix below; without it, the second read's
value visibly incorporated the first read's contribribution. \textbf{Fix:} fully clear the
per-read counters object at the start of every read, and pre-register (at zero) every counter
name the function's call graph touches (the complete, corrected list is: every name in the
reference source's own partial list, \emph{plus} the seven listed above).

\section{D16: Unconditional dereference of an empty optional on unmapped reads (real bug, UB)}
\label{sec:d16}
\textbf{Location:} \code{map\_read}, \texttt{src/shmap.h} (\S\ref{sec:map-read-output}).
\textbf{Description:} \code{best->set\_global\_stats(...)} and \code{best->print\_paf(...)} are
called unconditionally, textually after the point where the code has already branched on whether
\code{best} holds a value or not --- for every unmapped read, this dereferences an empty
\code{std::optional<Mapping>}, which is undefined behavior in C++ (it does not reliably crash;
depending on standard-library internals it may silently read whatever a default-constructed-
looking object's memory layout happens to contain). \textbf{Trigger:} any read that does not map
(a normal, expected, frequently-occurring outcome --- the surrounding code clearly anticipates it,
since a dedicated \code{unmapped\_out} stream and \code{lost\_on\_seeding}/\code{lost\_on\_pruning}
counters all exist specifically to track this case). \textbf{Fix:} guard both calls behind the
same \code{if (best)} check already present a few lines above, and define an explicit,
well-formed output for the unmapped case instead --- this document's companion implementation
writes a minimal 12-column PAF-shaped record (\S\ref{sec:unmapped-record}) rather than silently
omitting unmapped reads from output entirely, on the reasoning that the reference source's own
\code{unmapped\_out} stream was clearly intended to receive \emph{some} record per unmapped read,
even though the bug above meant it never cleanly did.

\section{D17: \code{gt\_C\_r\_lmax} is declared, printed, but never assigned}
\label{sec:d17}
\textbf{Location:} \code{AnalyseSimulatedReads} constructor, \texttt{src/analyse\_simulated.h}
(\S\ref{sec:ground-truth}). \textbf{Description:} distinct from D8 (which is a dead
\emph{parameter}) --- this is a genuinely missing assignment. The constructor computes and stores
\code{gt\_C\_l\_lmax} but has no corresponding \code{gt\_C\_r\_lmax = ...} line at all, despite
both fields being declared and both being unconditionally printed in \emph{both} the \code{-v 2}
PAF tags and the \code{.paul.tsv} output. \textbf{Impact:} the \code{gt\_C\_r\_lmax:f:} tag / TSV
column always shows the default \code{Mapping}'s sentinel score, $-1.000$, for every read, in
every run. \textbf{Recommendation:} since this diagnostic path has zero automated test coverage
to indicate the intended value, a from-scratch implementation may either leave the field
unassigned (matching shipped behavior exactly) or assign it the obviously-analogous computation
(mirroring \code{gt\_C\_l\_lmax} but for the right-side ground-truth bucket) --- both are
defensible; this document's companion implementation chose to leave it unassigned, to avoid
inventing behavior for a code path with no oracle to validate against.

\section{D18: PAF columns 10/11 always repeat the query-length column}
\label{sec:d18}
Already described in full in \S\ref{ch:output}; included here for completeness of this catalog.
Marked with the reference source's own \code{// TODO: fix} comment.

\section{D19: \code{Timers::secs}'s exception handling can never actually fire}
\label{sec:d19}
\textbf{Location:} \texttt{src/utils.h}. \textbf{Description:} \code{Timers::secs(name)} wraps
\code{timers\_.at(name).secs()} in a \code{try \{ ... \} catch (const std::runtime\_error\&)}
intended to gracefully handle a missing timer name. However, \code{std::unordered\_map::at} throws
\code{std::out\_of\_range} on a missing key, which does \emph{not} derive from
\code{std::runtime\_error} in the C++ standard exception hierarchy (both derive independently from
\code{std::exception}) --- so this catch clause can never actually trigger for the failure mode it
was written for. In practice this is masked because every timer name actually queried has always
been previously \code{.start()}'d (which auto-vivifies the entry), so the missing-key path is
never exercised by the shipped tool. \textbf{Recommendation:} a from-scratch implementation should
simply handle the ``timer not found'' case directly (e.g.\ return $0.0$ and log, which was
evidently the original intent) rather than reproduce a try/catch that provides no actual safety
net.

\section{D20: \code{ParsedQueryId::parse} catches a parse error only to rethrow it}
\label{sec:d20}
\textbf{Location:} \texttt{src/utils.h}. \textbf{Description:} when a ground-truth-encoded read
name's numeric fields (start/end position) fail to parse as integers, the function catches the
resulting exception, logs an error message, sets a (about-to-be-discarded) validity flag to false,
and then rethrows the same exception unconditionally --- so the net behavior is identical to not
catching it at all, just with an extra log line first. \textbf{Distinction from the two graceful
cases the same function handles} (no \code{!}-delimiter present at all; wrong number of
\code{!}-delimited fields), both of which \emph{do} return a clean ``not a ground-truth name''
result: a malformed \emph{numeric} field within an otherwise correctly-shaped name is treated as a
hard failure, not a graceful ``not ground-truth-encoded'' result. \textbf{Recommendation:}
this is reasonable, defensible behavior for a diagnostic feature that assumes well-formed
simulated-read names (only used in the opt-in \code{-v 2} mode); a from-scratch implementation
should simply let the parse error propagate (panic/exception) without the pointless
catch-log-rethrow wrapper, since the wrapper adds a log message but changes nothing else.
